\documentclass[10pt,a4paper]{article}
\usepackage[utf8]{inputenc}
\usepackage[vietnamese]{babel}
\usepackage{amsmath,amssymb,amsfonts}
\usepackage{geometry}
\geometry{a4paper, left=10mm, right=10mm, top=10mm, bottom=10mm}
\usepackage{tikz}
\usetikzlibrary{shapes,arrows.meta,positioning,calc}
\usepackage[most]{tcolorbox}
\usepackage{array,booktabs,multicol}
\usepackage{fontawesome}
\usepackage{helvet}
\renewcommand{\familydefault}{\sfdefault}
\usepackage{parskip}

\definecolor{primary}{RGB}{112, 48, 160}
\definecolor{teal}{RGB}{0, 128, 128}
\definecolor{orange}{RGB}{230, 108, 44}
\definecolor{green}{RGB}{46, 139, 87}
\definecolor{rose}{RGB}{192, 0, 0}

\pagestyle{empty}

\begin{document}

% --- HEADER BANNER ---
\begin{tcolorbox}[
  enhanced,
  colback=violet!8!white,
  colframe=violet!80!black,
  arc=2.5mm, boxrule=1pt,
  left=3mm, right=3mm, top=2mm, bottom=2mm
]
  \begin{tabular*}{\textwidth}{@{\extracolsep{\fill}} lr}
    {\small\bfseries\color{violet!85!black}\faBook\ SỔ TAY KIẾN THỨC TOÁN 2} & 
    {\small\bfseries\color{teal!85!black}\faGraduationCap\ CHƯƠNG TRÌNH GDPT 2018}
  \end{tabular*}
  \vspace{-1mm}
  \begin{center}
    \rule{0.98\textwidth}{0.6pt}
  \end{center}
  \vspace{-2mm}
  \centering
  {\LARGE\bfseries\color{violet!90!black}\faBookmark\ BÀI 62: PHÉP TRỪ (CÓ NHỚ) TRONG PHẠM VI 1000}\\[1mm]
  {\itshape\footnotesize\color{black!75} Kĩ thuật mượn và trả nhớ sang hàng chục hoặc hàng trăm khi trừ}
\end{tcolorbox}

\vspace{-2mm}

\noindent
\begin{minipage}[t]{0.49\textwidth}\vspace{0pt}
\begin{tcolorbox}[
  enhanced,
  colback=violet!4!white,
  colframe=violet!75!black,
  arc=2mm, boxrule=0.8pt,
  left=2mm, right=2mm, top=1.8mm, bottom=1.8mm,
  title={\footnotesize\bfseries\color{white}\faKey\ 1. Quy tắc trừ có nhớ}
]
{\footnotesize
\begin{itemize}\setlength{\itemsep}{1pt}\setlength{\leftskip}{-2.5mm}
  \item Khi trừ một hàng mà chữ số của số bị trừ nhỏ hơn chữ số của số trừ (không trừ được):
  \begin{itemize}\setlength{\itemsep}{0.5pt}\setlength{\leftskip}{-2mm}
    \item Ta \textbf{mượn 1} ở hàng liền kề bên trái (thành $10$ đơn vị của hàng hiện tại).
    \item Thực hiện trừ, rồi \textbf{trả 1} (thêm 1) vào chữ số hàng tiếp theo của số trừ.
  \end{itemize}
  \item \textbf{Khẩu quyết cốt lõi:} \textit{"Mượn ở đâu thì phải trả vào số trừ ở hàng bên cạnh"}.
  \item Trong chương trình Toán lớp 2, trừ có nhớ \textbf{không quá 1 lần}.
\end{itemize}
}
\end{tcolorbox}
\end{minipage}
\hfill
\begin{minipage}[t]{0.49\textwidth}\vspace{0pt}
\begin{tcolorbox}[
  enhanced,
  colback=teal!4!white,
  colframe=teal!75!black,
  arc=2mm, boxrule=0.8pt,
  left=2mm, right=2mm, top=1.8mm, bottom=1.8mm,
  title={\footnotesize\bfseries\color{white}\faSliders\ 2. Hai trường hợp có nhớ điển hình}
]
{\footnotesize
\begin{itemize}\setlength{\itemsep}{1pt}\setlength{\leftskip}{-2.5mm}
  \item \textbf{Trường hợp 1: Nhớ sang hàng chục}\\
  Ví dụ: $542 - 217$:
  \begin{itemize}\setlength{\itemsep}{0.5pt}\setlength{\leftskip}{-2mm}
    \item $2$ không trừ được $7$, lấy $12 - 7 = \mathbf{5}$, viết $5$ \textbf{nhớ $1$}.
    \item $1$ thêm $1$ là $2$; $4 - 2 = \mathbf{2}$, viết $2$.
    \item $5 - 2 = \mathbf{3}$, viết $3$.
    \item Kết quả: $\mathbf{542 - 217 = 325}$.
  \end{itemize}
  \item \textbf{Trường hợp 2: Nhớ sang hàng trăm}\\
  Ví dụ: $635 - 281$:
  \begin{itemize}\setlength{\itemsep}{0.5pt}\setlength{\leftskip}{-2mm}
    \item $5 - 1 = \mathbf{4}$, viết $4$.
    \item $3$ không trừ được $8$, lấy $13 - 8 = \mathbf{5}$, viết $5$ \textbf{nhớ $1$}.
    \item $2$ thêm $1$ là $3$; $6 - 3 = \mathbf{3}$, viết $3$.
    \item Kết quả: $\mathbf{635 - 281 = 354}$.
  \end{itemize}
\end{itemize}
}
\end{tcolorbox}
\end{minipage}

\vspace{-1.5mm}
\begin{tcolorbox}[
  enhanced,
  colback=orange!4!white,
  colframe=orange!80!black,
  arc=2mm, boxrule=0.8pt,
  left=2mm, right=2mm, top=1.5mm, bottom=1.5mm,
  title={\footnotesize\bfseries\color{white}\faSitemap\ 3. Sơ đồ quy trình mượn và trả 1 khi thực hiện trừ có nhớ}
]
\centering
\begin{tikzpicture}[x=1.65cm, y=0.75cm]
  \tikzset{arr/.style={-{Stealth[scale=0.8]}}, every node/.append style={align=center}}
  \node[rectangle, rounded corners=2pt, draw=violet!80!black, fill=violet!15, font=\scriptsize\bfseries] (m) at (0,0) {HÀNG HIỆN TẠI\\Không trừ được\\$\implies$ Lấy $10 + \dots$};
  \node[rectangle, rounded corners=2pt, draw=teal!80!black, fill=teal!15, font=\scriptsize\bfseries] (t) at (2.5,0) {TRỪ VÀ NHỚ\\Ghi kết quả\\Nhớ $1$};
  \node[rectangle, rounded corners=2pt, draw=green!70!black, fill=green!15, font=\scriptsize\bfseries] (r) at (5.0,0) {HÀNG BÊN CẠNH\\Thêm $1$ vào số trừ\\rồi mới thực hiện trừ};

  \draw[arr, thick] (m) -- (t);
  \draw[arr, thick, red!85!black] (t) -- node[above, font=\tiny\bfseries] {Trả 1} (r);
\end{tikzpicture}
\end{tcolorbox}

\vspace{-1.5mm}
\noindent
\begin{minipage}[t]{0.49\textwidth}\vspace{0pt}
\begin{tcolorbox}[
  enhanced,
  colback=green!4!white,
  colframe=green!60!black,
  arc=2mm, boxrule=0.8pt,
  left=2mm, right=2mm, top=1.8mm, bottom=1.8mm,
  title={\footnotesize\bfseries\color{white}\faCheckCircle\ 4. Ví dụ thực hành đặt tính}
]
{\footnotesize
\textbf{Đặt tính rồi tính:}\\
\begin{tabular}{r@{\,}c@{\,}lcr@{\,}c@{\,}l}
  & $653$ & & & & $718$ & \\
$-$ & $328$ & & &$-$ & $462$ & \\
\cline{1-3}\cline{5-7}
  & $\mathbf{325}$ & & & & $\mathbf{256}$ &
\end{tabular}
\vspace{1mm}\\
- $3$ không trừ được $8$, $13 - 8 = 5$, nhớ $1$; $2$ thêm $1$ là $3$; $5 - 3 = \mathbf{2}$; $6 - 3 = \mathbf{3}$.\\
- $8 - 2 = \mathbf{6}$; $1$ không trừ được $6$, $11 - 6 = 5$, nhớ $1$; $4$ thêm $1$ là $5$; $7 - 5 = \mathbf{2}$.
}
\end{tcolorbox}
\end{minipage}
\hfill
\begin{minipage}[t]{0.49\textwidth}\vspace{0pt}
\begin{tcolorbox}[
  enhanced,
  colback=rose!4!white,
  colframe=rose!70!black,
  arc=2mm, boxrule=0.8pt,
  left=2mm, right=2mm, top=1.8mm, bottom=1.8mm,
  title={\footnotesize\bfseries\color{white}\faExclamationTriangle\ 5. Lỗi quên trả nhớ tai hại}
]
{\footnotesize
\begin{itemize}\setlength{\itemsep}{1.5pt}\setlength{\leftskip}{-2.5mm}
  \item \textbf{Lỗi thường gặp nhất:} Mượn $1$ ở hàng đơn vị nhưng sang hàng chục quên không thêm $1$ vào số trừ $\implies$ kết quả bị sai lớn hơn $10$ đơn vị.
  \item \textbf{Lỗi lấy số dưới trừ số trên:} Khi thấy $2$ và $7$ ở hàng đơn vị, học sinh tự ý lấy $7 - 2 = 5$ là \textbf{hoàn toàn sai}!
  \item Phải lấy $12 - 7 = 5$ (mượn 1 chục).
\end{itemize}
}
\end{tcolorbox}
\end{minipage}


\end{document}
